%% ma: L12-B01-0039
%% lop: 12
%% chuong: 1
%% bai: 1
%% dang: 12.1.25
%% loai: TLN
%% muc: VD
%% dap_an: "1"
%% nguon: "HK1-DeOn_1_1.pdf (M10, Đề ôn tập học kì 1 Toán 12, Đề số 2), Phần III Câu 1"
%% kien_thuc: "Bài 1 (tính đơn điệu và cực trị), Lớp 11 Bài 32 (đạo hàm của hàm hợp)"
%% trang_thai: cho-duyet
%% ly_do: "Dạng mới đề xuất 12.1.25"
\begin{ex}
Cho hàm số $y=f(x)$ có đạo hàm $f'(x)=(-x+2)(x+1)$ với mọi $x\in\mathbb R$. Hàm số $f(x-3)$ đồng biến trên bao nhiêu khoảng?
\shortans{1}
\loigiai{$g(x)=f(x-3)$ có $g'(x)=f'(x-3)=(5-x)(x-2)$. $g'(x)>0\Leftrightarrow2<x<5$. Hàm số đồng biến trên đúng $1$ khoảng $(2;5)$.}
\end{ex}
